\subsection{Class-Term Extension}
\label{sec:technical-class-builders}

The primitive language of NBG contains no terms other than variables.
Nevertheless, the extended syntax introduced in this text permits
expressions of the form
\[
\{\xi\mid\varphi\}
\]
to be used as though they were terms denoting classes. We now specify
this extension and establish the formal basis for its use.

The intended role of such an expression is to represent the class
whose members are precisely the sets satisfying \(\varphi\). This
requires two distinct justifications. First, every class-builder
expression must represent a uniquely determined class. Second, every
complete formula containing class-builder expressions must admit
elimination in favor of a formula of the primitive language.

There is a further requirement arising from the predicative restriction
on Class Formation. A formula containing class-builder expressions may
itself occur as the defining condition of another class-builder
expression. We must therefore establish that the introduction of
class-builder notation does not permit the formation of classes whose
existence is not justified by the original Class Formation Schema.

These requirements will be addressed through a construction organized
by the complexity of class-builder expressions. We begin by specifying
the expressions involved and the procedures by which they are connected
to the existing language.

\subsubsection{Class-Term Syntax and Contextual Elimination}

\nbgpara{1}{General Class Terms}

We extend the syntax by admitting a new category of expressions, called
\emph{extended terms}.

\begin{definition}[General Class Terms]
\label{def:general-class-terms}
Every variable is an extended term.

If \(\varphi\) is a predicative extended formula and \(\xi\) is a
variable, then
\[
\{\xi\mid\varphi\}
\]
is a \emph{general class term}, provided that every class term occurring
in \(\varphi\) has already been admitted at an earlier stage of the
construction.

Every general class term is an extended term.

If \(S\) and \(T\) are extended terms, then
\[
S\in T
\qquad\text{and}\qquad
S=T
\]
are atomic extended formulas.
\end{definition}

The ordinary formation rules for formulas extend to these atomic
extended formulas. Thus, extended formulas containing class terms may
be combined using the logical connectives and quantifiers, as well as
previously admitted defined notation.

The variable \(\xi\) in
\[
\{\xi\mid\varphi\}
\]
binds every occurrence of \(\xi\) free in \(\varphi\). The free
variables of the resulting class term are therefore precisely the
variables occurring free in \(\varphi\), other than \(\xi\).

The conventions concerning free and bound occurrences, alphabetic
variation, and substitution of variables for variables extend to
class terms in the corresponding recursive manner. In particular,
the bound variable of a class term may be renamed alphabetically
to avoid capture. The compatibility of these operations with
primitive expansion will be established below.

The formation clauses do not yet license the substitution of arbitrary
extended terms for variables. This more general operation will be
justified separately, once the eliminability of class terms has been
established.

\nbgpara{2}{Bounded Class Terms}

We also admit a bounded form of class-builder notation.

\begin{definition}[Bounded Class Terms]
\label{def:bounded-class-terms}
If \(\xi\) and \(\eta\) are distinct variables and \(\varphi\) is a
predicative extended formula whose class terms, if any, have already
been admitted at earlier stages, then
\[
\{\xi\in\eta\mid\varphi\}
\]
is a \emph{bounded class term}.

Bounded class terms are defined by the expansion
\[
\{\xi\in\eta\mid\varphi\}
\defequiv
\{\xi\mid(\xi\in\eta\land\varphi)\}.
\]
\end{definition}

The variable \(\xi\) binds its free occurrences in \(\varphi\),
but does not bind \(\eta\). Thus, the free variables of
\[
\{\xi\in\eta\mid\varphi\}
\]
are the variables free in \(\varphi\) other than \(\xi\), together
with \(\eta\).

Bounded class terms introduce no independent class-forming operation.
Every occurrence is understood through its prescribed expansion into
a general class term. In particular, the existence and uniqueness
of the class represented by a bounded term will follow from the
corresponding results for general class terms.

The restriction to variable bounds is provisional. Once substitution
of extended terms has been justified, bounded class-builder notation
will be extended to arbitrary admitted extended terms as bounds.

\nbgpara{3}{Builder Rank}

Class terms may occur within the defining conditions of further class
terms. To organize their introduction without circularity, we assign
a \emph{builder rank} to extended expressions.

\begin{definition}[Builder Rank]
\label{def:builder-rank}
Variables have builder rank \(0\). An extended formula containing no
class term also has builder rank \(0\).

If
\[
T=\{\xi\mid\varphi\},
\]
then
\[
\operatorname{br}(T)
=
\operatorname{br}(\varphi)+1.
\]

A bounded class term has the builder rank of the general class term
to which it is expanded. Thus,
\[
\operatorname{br}
\bigl(
\{\xi\in\eta\mid\varphi\}
\bigr)
=
\operatorname{br}(\varphi)+1.
\]

The builder rank of any other extended expression containing class
terms is the greatest builder rank of the class terms occurring
within it.
\end{definition}

The formation of class terms is understood inductively in builder
rank.

At rank \(0\), the extended terms are the variables, and the
predicative extended formulas are those admitted before the
introduction of class-builder notation.

Suppose that class terms of rank at most \(n\), together with the
predicative extended formulas containing them, have been admitted.
A new class term
\[
\{\xi\mid\varphi\}
\]
may then be introduced if \(\varphi\) is predicative according to
the reductions already justified at these earlier ranks. Its rank
is
\[
\operatorname{br}(\varphi)+1.
\]

Consequently, a class term of rank \(n+1\) has a defining formula
containing only class terms of rank at most \(n\).

Predicativity and builder rank impose different requirements.
Predicativity determines which defining conditions may be used
in applications of Class Formation. Builder rank orders the
justification of nested class terms, ensuring that the admissibility
of a new term depends only upon results already established for
terms of lower rank.

The fact that this construction can be carried out through every
finite rank will be established by the Class-Term Extension
Metatheorem.

\nbgpara{4}{Canonical Membership Conditions}

We associate with each class term a formula expressing the condition
under which a variable belongs to the class represented by that term.

\begin{definition}[Canonical Membership Condition]
\label{def:technical-canonical-membership-condition}
Let
\[
T=\{\xi\mid\varphi\}
\]
be a general class term, and let \(\zeta\) be any variable.
Choose an alphabetic variant \(\varphi^\ast\) of \(\varphi\)
in which \(\zeta\) is free for \(\xi\). The
\emph{canonical membership condition} of \(T\), with
argument \(\zeta\), is
\[
\operatorname{Cond}_T(\zeta)
\defequiv
\Set(\zeta)\land\varphi^\ast[\zeta/\xi].
\]
\end{definition}

The sethood condition is essential. Under the intended interpretation
of NBG, every member of a class is a set. Thus, even if a proper class
satisfies the defining formula \(\varphi\), it cannot thereby become a
member of the class represented by \(T\).

For a bounded class term
\[
T=\{\xi\in\eta\mid\varphi\},
\]
the canonical membership condition is determined through its
expansion into a general class term:
\[
\operatorname{Cond}_T(\zeta)
\defequiv
\Set(\zeta)\land
\bigl(
\zeta\in\eta\land\varphi[\zeta/\xi]
\bigr).
\]
Since membership in \(\eta\) already implies sethood, this condition
is provably equivalent to
\[
\zeta\in\eta\land\varphi[\zeta/\xi].
\]
The resulting simplification will be used in the Membership Law for
bounded class terms.

The notation
\[
\operatorname{Cond}_T(\zeta)
\]
is metalinguistic. The expression ``\(\operatorname{Cond}\)''
is not a symbol of the primitive language or of the extended
syntax. It serves only to abbreviate the indicated extended formula.

Such a variant can always be obtained by alphabetically
renaming bound variables within \(\varphi\) as necessary
to avoid capture. Different legitimate choices yield
alphabetic variants of the resulting condition.

If \(T\) has builder rank \(n+1\), then
\(\operatorname{Cond}_T(\zeta)\) contains only class terms of rank
at most \(n\). This rank-lowering property will be fundamental
to the subsequent construction.

\nbgpara{5}{Realizations}

A class term is connected to the existing language by specifying
what it means for a variable to represent the class determined
by its canonical membership condition.

\begin{definition}[Realization of a Class Term]
\label{def:technical-realization-of-class-term}
Let \(T\) be a general class term, and let \(\eta\) be a variable.
Choose a variable \(\xi\) distinct from \(\eta\) and not occurring
in \(T\). Then
\[
\Real_T(\eta)
\defequiv
\forall\xi
\bigl(
\xi\in\eta
\leftrightarrow
\operatorname{Cond}_T(\xi)
\bigr).
\]
\end{definition}

The formula
\[
\Real_T(\eta)
\]
expresses that \(\eta\) has precisely the members specified by
the canonical membership condition of \(T\). When this formula
holds, we say that \(\eta\) \emph{realizes} \(T\), or that
\(\eta\) is a \emph{realization} of \(T\).

The notation \(\Real_T(\eta)\) is metalinguistic, just like
\(\operatorname{Cond}_T(\xi)\). It abbreviates an extended formula
and introduces no additional symbol into the primitive language.
Different legitimate choices of the bound variable \(\xi\)
yield alphabetic variants.

For bounded class terms, realization is defined through the
prescribed expansion into a general class term. Thus, if
\[
T=\{\xi\in\eta\mid\varphi\},
\]
then
\[
\Real_T(\zeta)
\]
abbreviates
\[
\Real_{\{\xi\mid(\xi\in\eta\land\varphi)\}}(\zeta).
\]

The definition does not itself assert that a realization exists.
In particular, an assertion of the form
\[
\exists\eta\,\Real_T(\eta)
\]
is intended to establish the existence of a class satisfying the
membership condition determined by the fixed parameters of \(T\).
For this purpose, \(\eta\) must not occur free in \(T\), since
otherwise the existential quantifier would also bind occurrences
belonging to those parameters.

The Class-Term Extension Metatheorem will establish that every
admitted class term has a realization, and that any two realizations
of the same class term are equal.

\nbgpara{6}{Contextual Elimination}

The primitive language contains no non-variable terms. A class
term therefore cannot itself be replaced by a term of the primitive
language. Instead, it must be eliminated within the complete
formula in which it occurs.

The elimination procedure uses fresh variables to represent
the classes specified by class terms.

\begin{definition}[Contextual Elimination]
\label{def:contextual-elimination}
Before the procedure is applied, every bounded class term is
expanded into its corresponding general class term.

Let \(\varphi\) be an extended formula containing at least one
class-term occurrence. Select the leftmost outermost such
occurrence, say \(T\), where an occurrence is \emph{outermost}
iff it does not lie within the defining formula of another
class-term occurrence.

The selected occurrence belongs to a unique atomic extended
subformula \(\theta\) of \(\varphi\).

Choose a variable \(\eta\) that does not occur in \(\varphi\),
and let \(\theta'\) be obtained from \(\theta\) by replacing
only the selected occurrence of \(T\) by \(\eta\).

Let \(\varphi'\) be obtained from \(\varphi\) by replacing
the selected occurrence of \(\theta\) with
\[
\exists\eta
\bigl(
\Real_T(\eta)\land\theta'
\bigr).
\]

A single step of \emph{contextual elimination} prescribes
\[
\varphi
\defequiv
\varphi'.
\]
The surrounding logical context is otherwise unchanged.
\end{definition}

\begin{example}
An occurrence of \(T \in A\)
is eliminated as
\[
\exists\eta
\bigl(
\Real_T(\eta)\land\eta\in A
\bigr).
\]
\end{example}

\begin{example}
An occurrence of \(A \in T\)
is eliminated as
\[
\exists\eta
\bigl(
\Real_T(\eta)\land A\in\eta
\bigr).
\]
\end{example}

\begin{example}
An occurrence of \(A=T\)
is eliminated as
\[
\exists\eta
\bigl(
\Real_T(\eta)\land A=\eta
\bigr).
\]
\end{example}

Only the selected occurrence of a class term is replaced.
If the same class term occurs more than once, its occurrences
are eliminated separately, using fresh realization variables
for each occurrence.

The procedure is repeated until no class-term occurrence
remains. Every step removes one class-builder occurrence:
the selected outer builder disappears, while the class terms
appearing in its canonical membership condition are precisely
those already occurring within its defining formula. Hence
the total number of class-builder occurrences decreases by one
at each step, and the procedure terminates after finitely
many steps.

The resulting expression contains no class terms. All remaining
defined notation can then be eliminated according to the
primitive-expansion clauses established previously.

\begin{definition}[Primitive Expansion of a Formula Containing Class Terms]
\label{def:primitive-expansion-with-class-terms}
A \emph{primitive expansion} of an extended formula containing
class terms is obtained by repeatedly applying contextual
elimination until no class terms remain, and then expanding
all remaining defined notation until a formula of the
primitive language is obtained.
\end{definition}

The procedure just specified applies to the extended formulas
generated by the class-term formation clauses, together with
the previously admitted logical and defined notation. The
application of earlier defined predicates and relations to
arbitrary class terms, where not yet licensed by their formation
clauses, will be justified separately through extended-term
substitution.

Although termination follows directly from the construction,
further properties require proof. In particular, we must show
that legitimate choices of fresh variables yield primitive
expansions differing only by alphabetic variation, and that
the procedure respects free variables and capture-avoiding
substitution.

These properties will be established by the Class-Term Extension
Metatheorem.

Contextual elimination must also be distinguished from
predicative reduction. The formula \(\Real_T(\eta)\) contains
an unrestricted universal quantifier, and contextual
elimination introduces an unrestricted existential quantifier.
Consequently, primitive expansion does not by itself establish
that a formula containing class terms may occur within a
predicative defining condition.

Separate predicative reductions will be required for that
purpose. Unlike contextual elimination, these reductions
will use only syntax admitted by Class Formation and will
replace atomic formulas containing class terms by equivalent
formulas involving terms of lower builder rank.

Together, contextual elimination and predicative reduction
will establish the two required aspects of class-builder
notation: its complete eliminability in favor of the
primitive language and its admissibility within the
predicative conditions used to form further classes.

\subsubsection{Class-Builder Metatheorem}

\nbgpara{1}{A Membership Abbreviation}

To formulate the predicative reductions uniformly, we introduce
an auxiliary metalinguistic abbreviation.

Let \(T\) be an extended term and let \(\zeta\) be a variable.
We write
\[
M_T(\zeta)
\]
for the following extended formula:
\[
M_T(\zeta)
\defequiv
\begin{cases}
\zeta\in T,
&
\text{if \(T\) is a variable},\\[2mm]
\operatorname{Cond}_T(\zeta),
&
\text{if \(T\) is a class term}.
\end{cases}
\]

In the second case, the bound variable of \(T\) is renamed
alphabetically whenever necessary to make the canonical membership
condition legitimate.

The abbreviation \(M_T(\zeta)\) expresses the membership condition
associated with \(T\), without presupposing that arbitrary extended
terms may already be substituted for variables.

Two properties will be useful.

First, if \(T\) is a class term of builder rank \(n+1\), then
\(M_T(\zeta)\) contains class terms only of rank at most \(n\).
If \(T\) is a variable, \(M_T(\zeta)\) contains no class terms.

Second, in either case,
\[
\vdash
M_T(\zeta)\to\Set(\zeta).
\]
For variable terms, this follows from Set-I. For general class
terms, it follows immediately from the definition of the
canonical membership condition. The bounded case follows through
its expansion into a general class term.

These observations concern only formulas containing terms already
admitted by the construction. They do not assert a general
substitution principle for extended terms.

\nbgpara{2}{The Class-Term Extension Metatheorem}

The definitions of the preceding paragraphs specify a stagewise
extension of the syntax, but they do not yet establish that the
construction can be continued through every finite builder rank.

In particular, the formation of a class term requires its defining
condition to be predicative. We must therefore show that formulas
containing previously admitted class terms have the predicative
reductions required for their use in further definitions.

At the same time, we must establish that each admitted term has a
unique realization and that contextual elimination yields
well-defined primitive expansions.

These properties will be proved together.

\begin{metatheorem}[Class-Term Extension]
\label{meta:class-term-extension}

The extension by general and bounded class terms is well defined
at every finite builder rank. More precisely:

\begin{enumerate}

\item Every admitted class term \(T\) has a realization.
If \(\eta\) does not occur free in \(T\), then
\[
\vdash
\exists\eta\,\Real_T(\eta).
\]

Any two realizations of \(T\) are equal. Thus, whenever the
displayed formulas are defined,
\[
\vdash
\bigl(
\Real_T(\eta)\land\Real_T(\zeta)
\bigr)
\to
\eta=\zeta.
\]

\item Repeated contextual elimination terminates and assigns every
extended formula containing class terms a primitive expansion,
determined up to alphabetic variation.

The expansion preserves the free variables of the original formula.
It is compatible with capture-avoiding substitution of variables
for variables, provided that freshly introduced bound variables
are renamed where necessary.

\item Every atomic extended formula containing a class term admits
a justified predicative reduction to an extended formula whose
class terms, if any, have strictly lower builder rank than the
greatest builder rank occurring in the original atomic formula.

The reduction introduces no new free variables.

\item Repeated application of these atomic reductions terminates.
Consequently, every formula constructed from the admitted
predicative syntax and class terms can be reduced to a
class-term-free predicative extended formula, and thence to a
core predicative formula by the previously justified reductions.

The resulting reduction is justified independently of any
application of Class Formation whose availability depends upon
that reduction.

\end{enumerate}
\end{metatheorem}

\begin{proof}
We proceed by induction on builder rank.

At rank \(0\), the extended terms are variables. The applicable
formation rules, primitive-expansion procedures, and predicative
reductions have already been established before the introduction
of class-builder notation. Thus, the required properties hold
at the initial stage.

Suppose that the construction has been justified through rank
\(n\). In particular, assume that the class terms of rank at
most \(n\) have unique realizations, that formulas containing
them admit primitive expansions with the stated syntactic
properties, and that the prescribed predicative reductions
have been justified for atomic formulas containing those terms.

We establish the corresponding results at rank \(n+1\).

\medskip

\noindent\textit{Existence and uniqueness of realizations.}

Let
\[
T=\{\xi\mid\varphi\}
\]
be an admitted general class term of rank \(n+1\).

By the formation condition for class terms, \(\varphi\) is
predicative according to the reductions already justified
at lower ranks. Moreover, every class term occurring in
\(\varphi\) has rank at most \(n\).

Choose variables \(\alpha\) and \(\eta\) not occurring in \(T\),
with \(\alpha\neq\eta\).

By the induction hypothesis, capture-avoiding substitution of
\(\alpha\) for \(\xi\) preserves the predicative admissibility
of \(\varphi\). Consequently,
\[
\operatorname{Cond}_T(\alpha)
\defequiv
\Set(\alpha)\land\varphi[\alpha/\xi]
\]
is predicative.

The conditions of Class Formation by Predicative Reduction
are therefore satisfied. Hence
\[
\vdash
\exists\eta\,\forall_S\alpha
\bigl(
\alpha\in\eta
\leftrightarrow
\operatorname{Cond}_T(\alpha)
\bigr).
\]

By the equivalence between restricted and unrestricted
membership specifications, this yields
\[
\vdash
\exists\eta\,\forall\alpha
\left(
\alpha\in\eta
\leftrightarrow
\bigl(
\Set(\alpha)\land
\operatorname{Cond}_T(\alpha)
\bigr)
\right).
\]

Since
\[
\vdash
\operatorname{Cond}_T(\alpha)
\to
\Set(\alpha),
\]
propositional reasoning simplifies the latter specification to
\[
\vdash
\exists\eta\,\forall\alpha
\bigl(
\alpha\in\eta
\leftrightarrow
\operatorname{Cond}_T(\alpha)
\bigr).
\]

By the definition of realization, this is
\[
\vdash
\exists\eta\,\Real_T(\eta).
\]

Thus, \(T\) has a realization.

For uniqueness, suppose that
\[
\Real_T(\eta)
\qquad\text{and}\qquad
\Real_T(\zeta).
\]
By their definitions,
\[
\forall\alpha
\bigl(
\alpha\in\eta
\leftrightarrow
\operatorname{Cond}_T(\alpha)
\bigr)
\]
and
\[
\forall\alpha
\bigl(
\alpha\in\zeta
\leftrightarrow
\operatorname{Cond}_T(\alpha)
\bigr).
\]
It follows that
\[
\forall\alpha
\bigl(
\alpha\in\eta
\leftrightarrow
\alpha\in\zeta
\bigr).
\]
Extensionality therefore gives
\[
\eta=\zeta.
\]

Existence and uniqueness for bounded class terms follow through
their prescribed expansions into general class terms.

\medskip

\noindent\textit{Contextual elimination and primitive expansion.}

Consider a formula whose class terms have rank at most \(n+1\).

A single step of contextual elimination replaces a selected
class-term occurrence \(T\) by a fresh variable \(\eta\) and
introduces the formula
\[
\Real_T(\eta).
\]

The selected outer builder is thereby removed. The class terms
occurring within \(\Real_T(\eta)\) are precisely those inherited
from the defining formula of \(T\), and hence have strictly
lower rank than \(T\).

As already observed in the definition of contextual elimination,
each step decreases the total number of class-builder occurrences
by one. The procedure therefore terminates after finitely many
steps.

When no class terms remain, all remaining defined notation can
be eliminated by the previously established expansion clauses.
The result is a formula of the primitive language.

The choice of a fresh realization variable does not affect this
result except by alphabetic variation. Indeed, two legitimate
choices introduce bound variables in corresponding positions,
and the resulting formulas are immediate alphabetic variants
after any necessary renaming of bound variables already
occurring in the surrounding expression.

Applying this observation successively to the finite elimination
procedure shows that the resulting primitive expansion is
determined up to alphabetic variation.

The procedure also preserves free variables. Replacing \(T\)
by \(\eta\) removes the free variables contributed by that
occurrence of \(T\), but the formula \(\Real_T(\eta)\)
reintroduces precisely those variables as free parameters.
The new variable \(\eta\) is bound by the introduced existential
quantifier. Thus, no free variable is gained or lost.

Compatibility with substitution of variables for variables
follows by induction on the formation of extended expressions.
In the atomic cases, contextual elimination commutes with
substitution after choosing its realization variables fresh
for both the original expression and the variable being
substituted. The logical connectives preserve this property
immediately. At quantifiers and class-builder binders, bound
variables are renamed alphabetically whenever necessary to
avoid capture.

The result is that substitution performed before contextual
elimination and substitution performed after contextual
elimination yield alphabetically variant primitive formulas,
provided that the fresh bound variables in the latter
expansion have first been chosen suitably.

The same argument establishes the corresponding compatibility
of alphabetic variation.

Hence the primitive-expansion procedure has the required
syntactic properties at rank \(n+1\).

\medskip

\noindent\textit{Membership and identity with variable arguments.}

We next establish two equivalences needed to justify the
predicative reductions.

Let \(T\) be a class term whose existence and uniqueness
have just been established, and let \(\zeta\) be a variable
for which its canonical membership condition is defined.

By contextual elimination,
\[
\zeta\in T
\defequiv
\exists\eta
\bigl(
\Real_T(\eta)\land\zeta\in\eta
\bigr),
\]
where \(\eta\) is fresh.

If \(\zeta\in T\), existential elimination supplies a
realization \(\eta\) with \(\zeta\in\eta\). The realization
condition gives
\[
\zeta\in\eta
\leftrightarrow
\operatorname{Cond}_T(\zeta),
\]
and hence \(\operatorname{Cond}_T(\zeta)\).

Conversely, suppose that
\[
\operatorname{Cond}_T(\zeta).
\]
Since a realization of \(T\) exists, existential elimination
supplies a variable \(\eta\) satisfying \(\Real_T(\eta)\).
The same membership equivalence yields
\[
\zeta\in\eta.
\]
Therefore,
\[
\exists\eta
\bigl(
\Real_T(\eta)\land\zeta\in\eta
\bigr),
\]
which is the contextual expansion of \(\zeta\in T\).

We have thus established
\[
\vdash
\zeta\in T
\leftrightarrow
\operatorname{Cond}_T(\zeta).
\]

Similarly, contextual elimination gives
\[
\zeta=T
\defequiv
\exists\eta
\bigl(
\Real_T(\eta)\land\zeta=\eta
\bigr).
\]

From this existential formula, identity elimination yields
\(\Real_T(\zeta)\). Conversely, from \(\Real_T(\zeta)\),
identity reflexivity and existential introduction yield the
displayed existential formula.

Consequently,
\[
\vdash
\zeta=T
\leftrightarrow
\Real_T(\zeta).
\]

Both arguments use only variable-based instances of the
previously established inference rules. They do not presuppose
the substitution of class terms for variables.

The resulting equivalences will be stated separately as the
fundamental Membership and Identity Laws below.

\medskip

\noindent\textit{Predicative reductions of atomic formulas.}

We now assign reductions to atomic extended formulas containing
class terms.

There are three cases. Throughout, variables introduced by a
reduction are chosen distinct from one another and absent from
the surrounding formula.

\smallskip

\noindent\emph{Case 1. Membership with a variable on the left.}

Let \(\zeta\) be a variable and \(T\) a class term.

The membership equivalence just established gives
\[
\vdash
\zeta\in T
\leftrightarrow
M_T(\zeta).
\]

We therefore prescribe
\[
\boxed{
\zeta\in T
\quad\rightsquigarrow\quad
M_T(\zeta).
}
\]

Since \(M_T(\zeta)\) contains only class terms of lower rank
than \(T\), this reduction strictly lowers builder rank.

\smallskip

\noindent\emph{Case 2. Membership with a class term on the left.}

Let \(S\) be a class term and \(T\) any extended term.

Choose fresh variables \(h\) and \(a\).

Contextual elimination of the occurrence of \(S\), followed
by the variable-argument membership equivalence for \(T\),
gives
\[
\vdash
S\in T
\leftrightarrow
\exists h
\bigl(
\Real_S(h)\land M_T(h)
\bigr).
\]

If \(T\) is a variable, the replacement of \(h\in T\) by
\(M_T(h)\) is simply an instance of the definition of \(M\).
If \(T\) is a class term, the replacement is justified by
the membership equivalence established above.

Now
\[
\vdash M_T(h)\to\Set(h).
\]
Furthermore, by the equivalence between restricted and
unrestricted membership specifications and the fact that
\(M_S(a)\) implies \(\Set(a)\), the realization condition
\[
\Real_S(h)
\]
is equivalent to
\[
\forall_S a
\bigl(
a\in h\leftrightarrow M_S(a)
\bigr).
\]

It follows that
\[
\vdash
S\in T
\leftrightarrow
\exists_S h
\left(
\forall_S a
\bigl(
a\in h\leftrightarrow M_S(a)
\bigr)
\land M_T(h)
\right).
\]

We therefore prescribe
\[
\boxed{
S\in T
\quad\rightsquigarrow\quad
\exists_S h
\left(
\forall_S a
\bigl(
a\in h\leftrightarrow M_S(a)
\bigr)
\land M_T(h)
\right).
}
\]

The reduction uses only restricted quantifiers. Every
class term occurring on the right belongs to the defining
formula of \(S\) or \(T\), and therefore has strictly lower
rank than the corresponding outer class term.

\smallskip

\noindent\emph{Case 3. Identity of extended terms.}

Let \(S\) and \(T\) be extended terms, at least one of which
is a class term.

Choose a fresh variable \(a\).

We claim that
\[
\vdash
S=T
\leftrightarrow
\forall_S a
\bigl(
M_S(a)\leftrightarrow M_T(a)
\bigr).
\]

To establish this, consider the realizations of any class
terms occurring in \(S\) and \(T\). Where one of these terms
is a variable, that variable itself is used.

By the variable-argument membership equivalence, the
members of each represented class are precisely the objects
satisfying its associated formula \(M_S(a)\) or \(M_T(a)\).

If the represented classes are equal, identity elimination
therefore gives agreement of their membership conditions.
Universal introduction, followed by restriction to sets,
establishes
\[
\forall_S a
\bigl(
M_S(a)\leftrightarrow M_T(a)
\bigr).
\]

Conversely, suppose that this restricted universal formula
holds.

Every member of any class is a set, and both \(M_S(a)\) and
\(M_T(a)\) imply \(\Set(a)\). Thus, the restricted equivalence
establishes agreement of the two membership conditions for
all objects.

Extensionality gives equality of the corresponding
realizations. Using the existence of realizations and
contextual elimination, we obtain \(S=T\).

Hence the claimed equivalence holds, and we prescribe
\[
\boxed{
S=T
\quad\rightsquigarrow\quad
\forall_S a
\bigl(
M_S(a)\leftrightarrow M_T(a)
\bigr).
}
\]

Again, every class term occurring on the right has strictly
lower rank than an outer class term occurring on the left.

\medskip

These three cases exhaust the atomic extended formulas
containing class terms. The only atomic relation symbols are
\(\in\) and \(=\), and the displayed cases cover every
possible placement of class terms in their arguments.

All three reductions are justified by derivable equivalences.
They introduce no free variables, since their newly
introduced variables are bound and chosen fresh for the
surrounding formula.

They also introduce no unrestricted quantifiers. Their
right-hand sides are constructed using the atomic and
previously admitted predicative notation, logical
connectives, restricted quantifiers, and class terms of
lower rank.

\medskip

\noindent\textit{Termination and predicative admissibility.}

It remains to establish that the prescribed reductions
can be applied repeatedly without circularity.

Consider the multiset of builder ranks of all class-term
occurrences in a finite extended formula.

Every atomic reduction removes the outer class terms
occurring in the selected atomic formula and replaces them
with formulas containing only their lower-ranked constituent
terms. Consequently, each application strictly decreases
the multiset of builder ranks under the standard multiset
ordering.

Since this ordering is well founded on finite multisets
of natural numbers, repeated reduction terminates.

More concretely, one may reduce all atomic occurrences
containing terms of the greatest rank, then those
containing terms of the next lower rank, and so on.
No reduction introduces a term of the rank just eliminated
or of any higher rank.

When no class terms remain, the resulting formula belongs
to the previously established extended syntax.

If the original formula was constructed using only the
admitted predicative operations, its class-term-free
reduct is predicative according to the reductions already
justified before class-builder notation was introduced.
Repeated application of those reductions therefore yields
a core predicative formula.

The equivalence of the original formula and its reduct
follows by successive applications of the atomic
equivalences, together with the established congruence
principles for the logical connectives and quantifiers.

All variable choices are made fresh for the surrounding
formula, so no reduction introduces additional free
variables or captures existing free occurrences. The
prescribed order of reduction determines the result up
to alphabetic variation.

Finally, the construction is non-circular. At rank
\(n+1\), the existence of realizations was established
using only predicative reductions justified at ranks at
most \(n\). The membership and identity equivalences
were then derived from those realizations. Only afterward
were the new atomic reductions justified.

Thus, no application of Class Formation depends upon
a reduction whose justification presupposes that same
application.

The induction step is complete. By induction, the stated
properties hold at every finite builder rank.
\end{proof}


\subsubsection{Fundamental Laws and Derived Rules}

\nbgpara{1}{Fundamental Laws for General Class Terms}

The Class-Term Extension Metatheorem establishes that every
admitted class term has a unique realization and that formulas
containing class terms admit elimination into the primitive
language. We can now state explicitly the two fundamental laws
governing membership in and identity with class terms.

These laws are derived theorems. They are not additional clauses
in the definition of class-builder notation, nor do they
introduce new inference principles.

\begin{metatheorem}[Membership Law for General Class Terms]
\label{meta:membership-law-for-class-terms}
Let \(T\) be a general class term, and let \(\zeta\) be a variable
for which \(\operatorname{Cond}_T(\zeta)\) is defined. Then
\[
\vdash
\zeta\in T
\leftrightarrow
\operatorname{Cond}_T(\zeta).
\]
\end{metatheorem}

\begin{proof}
The equivalence was established in the proof of
Metatheorem~\ref{meta:class-term-extension}.

Recall that contextual elimination identifies
\(\zeta\in T\) with
\[
\exists\eta
\bigl(
\Real_T(\eta)\land\zeta\in\eta
\bigr),
\]
where \(\eta\) is fresh.

Every realization of \(T\) satisfies
\[
\zeta\in\eta
\leftrightarrow
\operatorname{Cond}_T(\zeta).
\]
The forward implication therefore follows immediately by
existential elimination.

For the converse, the existence of a realization of \(T\)
allows us to infer, from
\(\operatorname{Cond}_T(\zeta)\), that some realization
\(\eta\) contains \(\zeta\). Existential introduction then
gives the contextual expansion of \(\zeta\in T\).

The two implications establish the equivalence.
\end{proof}

For a general class term
\[
T=\{\xi\mid\varphi\},
\]
the law takes the concrete form
\[
\boxed{
\vdash
\zeta\in\{\xi\mid\varphi\}
\leftrightarrow
\bigl(
\Set(\zeta)\land\varphi[\zeta/\xi]
\bigr).
}
\]

The sethood condition is not dispensable. The formula
\(\varphi[\zeta/\xi]\) may hold even when \(\zeta\) is a
proper class, whereas membership in any class entails
sethood.

The corresponding law for identity characterizes when a
variable represents the class specified by a class term.

\begin{metatheorem}[Identity Law for General Class Terms]
\label{meta:identity-law-for-class-terms}
Let \(T\) be a general class term, and let \(\zeta\) be a
variable for which \(\Real_T(\zeta)\) is defined. Then
\[
\vdash
\zeta=T
\leftrightarrow
\Real_T(\zeta).
\]
\end{metatheorem}

\begin{proof}
This equivalence was likewise established in the proof of
Metatheorem~\ref{meta:class-term-extension}.

By contextual elimination,
\[
\zeta=T
\]
is definitionally equivalent to
\[
\exists\eta
\bigl(
\Real_T(\eta)\land\zeta=\eta
\bigr),
\]
where \(\eta\) is fresh.

From this formula, existential elimination and identity
elimination yield \(\Real_T(\zeta)\).

Conversely, from \(\Real_T(\zeta)\), identity reflexivity
gives
\[
\Real_T(\zeta)\land\zeta=\zeta.
\]
Existential introduction therefore yields the contextual
expansion of \(\zeta=T\).

The two implications establish the equivalence.
\end{proof}

Expanding the realization condition gives the concrete form
\[
\boxed{
\vdash
\zeta=\{\xi\mid\varphi\}
\leftrightarrow
\forall\alpha
\left(
\alpha\in\zeta
\leftrightarrow
\bigl(
\Set(\alpha)\land\varphi[\alpha/\xi]
\bigr)
\right),
}
\]
where \(\alpha\) is chosen fresh for the displayed formula.

The Membership Law and Identity Law express two different
aspects of class-builder notation. The former determines the
members of the class represented by a term. The latter
characterizes that class by its complete membership
specification.

Both laws are consequences of the original deductive system
together with Class Formation and Extensionality. Their use
does not enlarge the primitive language or the stock of
axioms.


\nbgpara{2}{Fundamental Laws for Bounded Class Terms}

The corresponding laws for bounded class terms follow from
their prescribed expansions into general class terms.

\begin{metatheorem}[Fundamental Laws for Bounded Class Terms]
\label{meta:fundamental-laws-for-bounded-class-terms}
Let
\[
T=\{\xi\in\eta\mid\varphi\}
\]
be an admitted bounded class term, where \(\xi\) and \(\eta\)
are distinct variables.

Whenever the indicated substitutions are defined,
\[
\vdash
\zeta\in T
\leftrightarrow
\bigl(
\zeta\in\eta
\land
\varphi[\zeta/\xi]
\bigr),
\]
and
\[
\vdash
\zeta=T
\leftrightarrow
\forall\alpha
\left(
\alpha\in\zeta
\leftrightarrow
\bigl(
\alpha\in\eta
\land
\varphi[\alpha/\xi]
\bigr)
\right),
\]
where \(\alpha\) is chosen fresh.
\end{metatheorem}

\begin{proof}
By Definition~\ref{def:bounded-class-terms},
\[
T
\defequiv
\{\xi\mid(\xi\in\eta\land\varphi)\}.
\]

Applying the Membership Law for general class terms gives
\[
\vdash
\zeta\in T
\leftrightarrow
\left(
\Set(\zeta)\land
\bigl(
\zeta\in\eta
\land
\varphi[\zeta/\xi]
\bigr)
\right).
\]

However, Set-I gives
\[
\vdash
\zeta\in\eta\to\Set(\zeta).
\]
Consequently, the sethood conjunct is redundant, and
propositional reasoning yields
\[
\vdash
\zeta\in T
\leftrightarrow
\bigl(
\zeta\in\eta
\land
\varphi[\zeta/\xi]
\bigr).
\]

For identity, the Identity Law for general class terms gives
\[
\vdash
\zeta=T
\leftrightarrow
\forall\alpha
\left(
\alpha\in\zeta
\leftrightarrow
\left(
\Set(\alpha)\land
\bigl(
\alpha\in\eta
\land
\varphi[\alpha/\xi]
\bigr)
\right)
\right).
\]

Since
\[
\vdash
\alpha\in\eta\to\Set(\alpha),
\]
the equivalence
\[
\vdash
\left(
\Set(\alpha)\land
\bigl(
\alpha\in\eta
\land
\varphi[\alpha/\xi]
\bigr)
\right)
\leftrightarrow
\bigl(
\alpha\in\eta
\land
\varphi[\alpha/\xi]
\bigr)
\]
follows by propositional reasoning.

Quantifier Congruence and the transitivity of provable
equivalence therefore yield
\[
\vdash
\zeta=T
\leftrightarrow
\forall\alpha
\left(
\alpha\in\zeta
\leftrightarrow
\bigl(
\alpha\in\eta
\land
\varphi[\alpha/\xi]
\bigr)
\right).
\]
\end{proof}

Thus, although bounded class-builder notation is defined by
expansion into general class-builder notation, its fundamental
laws admit simpler formulations. In particular, the Membership
Law for bounded class terms requires no explicit sethood
condition, since membership in the bounding class already
guarantees it.

The simplified membership conditions are provably equivalent
to the canonical membership conditions introduced earlier.
They are not additional definitional expansions.


\nbgpara{3}{Derived Rules for Class Terms}

The Membership and Identity Laws each establish a biconditional.
Their two directions may therefore be extracted as derived
inference rules.

\[
\frac{\zeta\in T}
{\operatorname{Cond}_T(\zeta)}
\quad
\mathrm{CB}{\in}\text{-E}
\qquad\qquad
\frac{\operatorname{Cond}_T(\zeta)}
{\zeta\in T}
\quad
\mathrm{CB}{\in}\text{-I}
\]

\[
\frac{\zeta=T}
{\Real_T(\zeta)}
\quad
\mathrm{CB}{=}\text{-E}
\qquad\qquad
\frac{\Real_T(\zeta)}
{\zeta=T}
\quad
\mathrm{CB}{=}\text{-I}
\]

Here, \(T\) is an admitted class term and \(\zeta\) is a
variable.

The membership rules are permitted whenever
\(\operatorname{Cond}_T(\zeta)\) is defined, and the identity
rules are permitted whenever \(\Real_T(\zeta)\) is defined.

The bound variable of a class term may be renamed
alphabetically whenever necessary to make the relevant
substitution legitimate. Variables introduced in the expansion
of a realization condition are chosen fresh for the surrounding
expression.

For bounded class terms, the simplified membership condition
established in the preceding paragraph may be used in place
of the canonical condition. Thus, for example, the rules
\[
\frac{
\zeta\in\{\xi\in\eta\mid\varphi\}
}{
\zeta\in\eta\land\varphi[\zeta/\xi]
}
\quad
\mathrm{CB}{\in}\text{-E}
\]
and
\[
\frac{
\zeta\in\eta\land\varphi[\zeta/\xi]
}{
\zeta\in\{\xi\in\eta\mid\varphi\}
}
\quad
\mathrm{CB}{\in}\text{-I}
\]
are admissible whenever the indicated substitution is defined.

Each application of a class-builder rule abbreviates an
application of the corresponding fundamental law, followed
by Biconditional Elimination. The rules therefore introduce
no new deductive power.

The rules just stated have variables as their explicit
arguments. The more general use of class terms as arguments,
as in
\[
S\in T
\]
with \(S\) itself a class term, has already been assigned a
primitive expansion and a predicative reduction by the
Class-Term Extension Metatheorem. Its treatment by the
ordinary inference rules nevertheless requires a further
admissibility result concerning substitution of extended
terms for variables.

Until that result has been established, the present rules
are understood in their variable-argument forms.


\nbgpara{4}{Class-Builder Congruence}

The preceding laws characterize class terms through their
membership conditions. Together with Extensionality, they
also provide a general criterion for equality between
class-builder expressions.

Two general class terms represent the same class precisely
when their defining conditions agree on all sets. Agreement
on proper classes is irrelevant, since proper classes cannot
be members of either class.

\begin{metatheorem}[Class-Builder Congruence]
\label{meta:class-builder-congruence}
Let \(\varphi\) and \(\psi\) be predicative extended formulas
such that
\[
\{\xi\mid\varphi\}
\qquad\text{and}\qquad
\{\xi\mid\psi\}
\]
are admitted class terms. Then
\[
\vdash
\{\xi\mid\varphi\}
=
\{\xi\mid\psi\}
\leftrightarrow
\forall_S\xi
\bigl(
\varphi\leftrightarrow\psi
\bigr).
\]
\end{metatheorem}

\begin{proof}
Let
\[
T=\{\xi\mid\varphi\}
\qquad\text{and}\qquad
U=\{\xi\mid\psi\}.
\]

Choose a variable \(\alpha\) distinct from \(\xi\) and not
occurring in \(T\) or \(U\).

By the atomic identity reduction established in
Metatheorem~\ref{meta:class-term-extension},
\[
\vdash
T=U
\leftrightarrow
\forall_S\alpha
\bigl(
M_T(\alpha)\leftrightarrow M_U(\alpha)
\bigr).
\]

By the definition of the auxiliary membership abbreviation,
\[
M_T(\alpha)
\defequiv
\Set(\alpha)\land\varphi[\alpha/\xi]
\]
and
\[
M_U(\alpha)
\defequiv
\Set(\alpha)\land\psi[\alpha/\xi].
\]

Consequently,
\[
\vdash
T=U
\leftrightarrow
\forall_S\alpha
\left(
\bigl(
\Set(\alpha)\land\varphi[\alpha/\xi]
\bigr)
\leftrightarrow
\bigl(
\Set(\alpha)\land\psi[\alpha/\xi]
\bigr)
\right).
\]

Within the scope of a restricted universal quantifier,
\(\Set(\alpha)\) may be assumed. Under this assumption,
propositional reasoning establishes
\[
\begin{aligned}
&
\Bigl(\bigl(
\Set(\alpha)\land\varphi[\alpha/\xi]
\bigr)
\leftrightarrow
\bigl(
\Set(\alpha)\land\psi[\alpha/\xi]
\bigr)\Bigr)\leftrightarrow
\bigl(
\varphi[\alpha/\xi]
\leftrightarrow
\psi[\alpha/\xi]
\bigr).
\end{aligned}
\]

By Restricted Quantifier Congruence, it follows that
\[
\vdash
T=U
\leftrightarrow
\forall_S\alpha
\bigl(
\varphi[\alpha/\xi]
\leftrightarrow
\psi[\alpha/\xi]
\bigr).
\]

Since \(\alpha\) was chosen fresh, the formula on the right
is an alphabetic variant of
\[
\forall_S\xi
\bigl(
\varphi\leftrightarrow\psi
\bigr).
\]

Replacement by alphabetic variants therefore gives
\[
\vdash
\{\xi\mid\varphi\}
=
\{\xi\mid\psi\}
\leftrightarrow
\forall_S\xi
\bigl(
\varphi\leftrightarrow\psi
\bigr),
\]
as required.
\end{proof}

The metatheorem strengthens the familiar principle that
equivalent defining conditions determine equal classes.
In particular, if
\[
\vdash
\forall\xi
\bigl(
\varphi\leftrightarrow\psi
\bigr),
\]
then Restricted Universal Elimination of the unrestricted
specification, together with Class-Builder Congruence, gives
\[
\vdash
\{\xi\mid\varphi\}
=
\{\xi\mid\psi\}.
\]

The unrestricted equivalence is sufficient but not
necessary. The two defining conditions need agree only
when their argument is a set.

The corresponding criterion for bounded class terms
follows immediately through their expansions into
general class terms. For example, if \(\eta\) is distinct
from \(\xi\) and both bounded terms are admitted, then
\[
\vdash
\begin{aligned}[t]
&
\{\xi\in\eta\mid\varphi\}
=
\{\xi\in\eta\mid\psi\}\leftrightarrow
\forall_S\xi
\bigl(
\xi\in\eta
\to
(\varphi\leftrightarrow\psi)
\bigr).
\end{aligned}
\]

Thus, defining conditions for bounded class terms need
agree only on the sets belonging to the bounding class.

Class-Builder Congruence will also be useful in justifying
the substitution of extended terms. In particular, it
allows us to establish equality between class terms whose
defining conditions have been transformed by provably
equivalent replacements.

Unlike the original rankwise construction, this congruence
result does not participate in the justification of
predicative reductions. It is a consequence of the
Class-Term Extension Metatheorem, rather than a prerequisite
for that theorem.





\subsubsection{Extended-Term Substitution and Admissibility}
\label{sec:extended-term-substitution}

The Class-Term Extension Metatheorem establishes that class terms
have unique realizations and that formulas containing them admit
primitive expansions. It also justifies their occurrence within
predicative defining conditions.

These results do not yet establish that class terms may be
substituted for variables in arbitrary formulas and derivations.
The primitive language contains only variables as terms, and
the inference rules of Gentzen ND were formulated accordingly.

We now establish that the extended syntax nevertheless admits
the ordinary operations associated with terms. These operations
will be justified through realizations and contextual elimination.
They introduce no additional primitive inference rules.

\nbgpara{1}{Representation by Variables}

The realization relation provides a means of replacing a class
term by a variable representing the same class. For the present
purpose, it is convenient to extend this method uniformly to
all extended terms, including variables.

\begin{definition}[Representation of Extended Terms]
\label{def:representation-of-extended-terms}
Let \(T\) be an extended term, and let \(\eta\) be a variable.
We introduce the following metalinguistic abbreviation:
\[
\operatorname{Rep}_T(\eta)
\defequiv
\begin{cases}
\eta=T,
&
\text{if \(T\) is a variable},\\[2mm]
\Real_T(\eta),
&
\text{if \(T\) is a class term}.
\end{cases}
\]
\end{definition}

Thus, a variable represents another variable precisely when
the two are identical. A variable represents a class term
precisely when it realizes that term.

This notation is auxiliary. In particular, it does not
introduce expressions of the form \(\Real_T(S)\) with
arbitrary extended terms as arguments.

The existence and uniqueness of representatives follow
immediately from the preceding development. If \(\eta\)
is fresh for \(T\), then
\[
\vdash
\exists\eta\,\operatorname{Rep}_T(\eta).
\]
Moreover, whenever the displayed formulas are defined,
\[
\vdash
\bigl(
\operatorname{Rep}_T(\eta)
\land
\operatorname{Rep}_T(\zeta)
\bigr)
\to
\eta=\zeta.
\]

For class terms, these are consequences of the Class-Term
Extension Metatheorem. For variables, they follow from
the primitive rules for identity.

The following observation will be useful.

\begin{metatheorem}[Representation of Atomic Formulas]
\label{meta:representation-of-atomic-formulas}
Let \(S\) and \(T\) be extended terms, and let \(\eta\)
and \(\zeta\) be distinct variables chosen fresh for
both terms. Then
\[
\left\{
\operatorname{Rep}_S(\eta),
\operatorname{Rep}_T(\zeta)
\right\}
\vdash
(S\in T\leftrightarrow\eta\in\zeta)
\]
and
\[
\left\{
\operatorname{Rep}_S(\eta),
\operatorname{Rep}_T(\zeta)
\right\}
\vdash
(S=T\leftrightarrow\eta=\zeta).
\]
\end{metatheorem}

\begin{proof}
If both terms are variables, the equivalences follow
directly from the primitive identity rules.

Otherwise, consider the contextual elimination of the
class terms occurring in the displayed atomic formulas.

The resulting formulas assert the existence of variables
realizing the respective class terms, together with the
corresponding primitive membership or identity relation.

Under the assumptions
\[
\operatorname{Rep}_S(\eta)
\qquad\text{and}\qquad
\operatorname{Rep}_T(\zeta),
\]
the uniqueness of representatives allows every
existentially introduced realization variable to be
identified with \(\eta\) or \(\zeta\), respectively.

The primitive identity-elimination rules then establish
the required atomic equivalences. Conversely, the
existence of the displayed representatives permits the
corresponding existential statements to be introduced.

Thus, the contextual expansions of the original atomic
formulas are equivalent to the primitive relations between
their representatives.
\end{proof}

Consequently, the primitive identity rules also justify
the replacement of extended terms by equal extended terms
within atomic formulas. This follows by choosing
representatives, applying the primitive rules to those
representatives, and eliminating the realization variables.

The generalization to arbitrary extended formulas will
be established below.


\nbgpara{2}{Defined Relations with Extended-Term Arguments}

Before introducing general extended-term substitution, we
extend the basic defined relations of NBG to arbitrary
extended terms.

For an extended term \(T\), we admit
\[
\Set(T)
\defequiv
\exists\eta(T\in\eta),
\]
where \(\eta\) is fresh for \(T\), and
\[
\PC(T)
\defequiv
\neg\Set(T).
\]

For extended terms \(S\) and \(T\), we admit
\[
S\subseteq T
\defequiv
\forall\xi
\bigl(
\xi\in S\to\xi\in T
\bigr),
\]
where \(\xi\) is fresh for both terms.

These definitions agree with the previously established
definitions when their arguments are variables. Since the
right-hand sides are extended formulas already admitted
by the class-term syntax, they also provide contextual
eliminations for the newly formed expressions.

Their predicative admissibility requires a separate
justification.

For this purpose, let
\[
\operatorname{Spec}_T(\eta)
\defequiv
\forall_S\xi
\bigl(
\xi\in\eta
\leftrightarrow
M_T(\xi)
\bigr),
\]
where \(\xi\) is fresh for \(T\) and \(\eta\).

By the restricted and unrestricted membership
specification equivalence, together with the fact that
\(M_T(\xi)\) implies \(\Set(\xi)\), we have
\[
\vdash
\operatorname{Spec}_T(\eta)
\leftrightarrow
\operatorname{Rep}_T(\eta).
\]

For class terms, this is the equivalence between
restricted membership specification and realization.
For variables, it follows additionally from Extensionality.

We can therefore establish the following reductions.

\begin{metatheorem}[Predicative Reduction of Basic Term Relations]
\label{meta:predicative-reduction-of-term-relations}
Let \(S\) and \(T\) be extended terms. With all newly
introduced variables chosen fresh, the following
equivalences hold:
\[
\vdash
\Set(T)
\leftrightarrow
\exists_S\eta\,
\operatorname{Spec}_T(\eta),
\]
and
\[
\vdash
S\subseteq T
\leftrightarrow
\forall_S\xi
\bigl(
M_S(\xi)\to M_T(\xi)
\bigr).
\]
\end{metatheorem}

\begin{proof}
For sethood, choose a fresh representative \(\eta\)
of \(T\).

By the representation equivalence for membership,
\(\Set(T)\) holds precisely when its representative
\(\eta\) is a set. Since every extended term has a
representative, we obtain
\[
\vdash
\Set(T)
\leftrightarrow
\exists\eta
\bigl(
\operatorname{Rep}_T(\eta)
\land\Set(\eta)
\bigr).
\]

Using the equivalence between representation and
\(\operatorname{Spec}_T(\eta)\), the right-hand side is
equivalent to
\[
\exists_S\eta\,
\operatorname{Spec}_T(\eta).
\]

For subclasshood, the defining formula is
\[
\forall\xi
\bigl(
\xi\in S\to\xi\in T
\bigr).
\]

By the variable-argument Membership Law, membership
in either extended term is equivalent to its
associated condition \(M\). Thus,
\[
\vdash
S\subseteq T
\leftrightarrow
\forall\xi
\bigl(
M_S(\xi)\to M_T(\xi)
\bigr).
\]

Since \(M_S(\xi)\) implies \(\Set(\xi)\), restricting
the universal quantifier to sets does not change
the truth of the implication. This gives the
required equivalence.
\end{proof}

We assign the right-hand sides of these equivalences
as predicative reductions of the corresponding
defined relations.

The reduction for \(\PC(T)\) is obtained by negating
the reduction for \(\Set(T)\). The reductions for
negated membership, identity, and subclasshood
relations are likewise obtained by negation.

Each reduction uses only the previously admitted
predicative syntax and canonical membership conditions
of lower builder rank. Consequently, these defined
relations may occur within the predicative conditions
of further class terms.

These assignments are required for the extensions
just introduced. Other defined notation may be
extended to class-term arguments in the same manner
when its formation and predicative-reduction clauses
have been justified.


\nbgpara{3}{Substitution of Extended Terms}

We can now extend the substitution operation beyond
the replacement of variables by variables.

\begin{definition}[Extended-Term Substitution]
\label{def:extended-term-substitution}
Let \(\varphi\) be an extended formula, let \(\xi\)
be a variable, and let \(T\) be an extended term.

The expression
\[
\varphi[T/\xi]
\]
denotes the result of replacing every free occurrence
of \(\xi\) in \(\varphi\) by \(T\), with bound variables
renamed alphabetically whenever necessary to prevent
capture of free variables occurring in \(T\).

The same operation is defined recursively for extended
terms.
\end{definition}

At quantifiers, substitution proceeds into the scope
of the quantifier unless that quantifier binds \(\xi\).
At class-builder expressions, it proceeds into the
defining formula unless the builder binds \(\xi\).

Whenever a binder would capture a free variable of
the substituted term, its bound variable is renamed
to a suitable fresh variable before substitution.

The operation is extended to defined notation in
accordance with its formation and expansion clauses.
In particular, the defined relations introduced in
the preceding paragraph are treated through their
prescribed expansions.

Unlike substitution of variables for variables,
extended-term substitution need not preserve builder
rank. Substitution of a class term into the defining
condition of another class term may increase the rank
of the resulting expression.

This does not present a difficulty for the stagewise
construction, provided that the resulting defining
condition remains predicative.

\begin{metatheorem}[Preservation of Predicative Admissibility]
\label{meta:predicativity-under-term-substitution}
Let \(\varphi\) be a predicative extended formula,
let \(\xi\) be a variable, and let \(T\) be an
admitted extended term.

Then the capture-avoiding substitution
\[
\varphi[T/\xi]
\]
is predicative, provided that all defined notation
occurring in \(\varphi\) is among the notation whose
extended-term formation and predicative-reduction
clauses have been justified.
\end{metatheorem}

\begin{proof}
Proceed by induction on the formation of the
predicative extended formula.

For atomic formulas, substitution yields atomic
extended formulas containing admitted extended
terms. The Class-Term Extension Metatheorem supplies
their justified rank-lowering reductions.

The defined relations treated in the preceding
paragraph have likewise been assigned justified
predicative reductions for arbitrary extended-term
arguments.

The logical connectives preserve predicative
admissibility. The same holds for restricted
quantification, after alphabetic renaming of
bound variables where necessary.

If the substitution occurs within the defining
formula of an existing class term, the induction
hypothesis establishes that the resulting defining
condition is predicative. The substituted class
term and all constituent terms have finite builder
rank, so the resulting builder is admitted at the
rank determined by its new defining formula.

Thus, every stage of the substitution produces
an admitted extended expression, and the resulting
formula is predicative.
\end{proof}

This theorem concerns substitution within already
admitted predicative syntax. It does not permit an
unrestricted quantifier over classes to occur in a
class-forming condition merely because that
quantifier has been concealed within defined
notation.


\nbgpara{4}{Realization Invariance}

The central justification of extended-term
substitution is that an extended formula cannot
distinguish a class term from a variable realizing
that term.

\begin{metatheorem}[Realization Invariance]
\label{meta:realization-invariance}
Let \(\varphi\) be an extended formula, let \(\xi\)
be a variable, and let \(T\) be an admitted class
term.

Choose a variable \(\eta\) that does not occur
in \(\varphi\) or \(T\), and assume that all
substitutions are capture-avoiding.

Then
\[
\{
\Real_T(\eta)
\}
\vdash
\varphi[T/\xi]
\leftrightarrow
\varphi[\eta/\xi].
\]

Consequently,
\[
\boxed{
\vdash
\varphi[T/\xi]
\leftrightarrow
\exists\eta
\bigl(
\Real_T(\eta)
\land
\varphi[\eta/\xi]
\bigr).
}
\]
\end{metatheorem}

\begin{proof}
We first establish the equivalence under the assumption
\(\Real_T(\eta)\).

The proof proceeds by simultaneous structural induction
on extended terms and extended formulas. For an extended
term \(U\), the corresponding induction statement is
\[
\{
\Real_T(\eta)
\}
\vdash
U[T/\xi]=U[\eta/\xi].
\]

For variables, this follows from identity reflexivity
and the Identity Law for \(T\). In particular,
\(\Real_T(\eta)\) entails both \(\eta=T\) and \(T=\eta\)
by their contextual expansions.

Suppose next that
\[
U=\{\zeta\mid\psi\}
\]
is a class term.

If \(\zeta=\xi\), the binder prevents substitution
into the occurrences of \(\xi\) that it binds, and
the required equality is immediate.

Otherwise, rename \(\zeta\) alphabetically so that
it is fresh for \(T\) and \(\eta\).

By the induction hypothesis applied to the defining
formula \(\psi\),
\[
\{
\Real_T(\eta)
\}
\vdash
\psi[T/\xi]
\leftrightarrow
\psi[\eta/\xi].
\]

Since \(\zeta\) does not occur free in the remaining
assumption \(\Real_T(\eta)\), restricted universal
introduction gives
\[
\{
\Real_T(\eta)
\}
\vdash
\forall_S\zeta
\bigl(
\psi[T/\xi]
\leftrightarrow
\psi[\eta/\xi]
\bigr).
\]

Class-Builder Congruence therefore yields
\[
\{
\Real_T(\eta)
\}
\vdash
\{\zeta\mid\psi[T/\xi]\}
=
\{\zeta\mid\psi[\eta/\xi]\}.
\]

The case of bounded class terms follows through
their prescribed expansion into general class terms.

Now consider atomic formulas. The term induction
establishes equality between the corresponding
arguments before and after substitution.

By Metatheorem~\ref{meta:representation-of-atomic-formulas},
these equalities may be transferred to their
primitive representatives. The primitive identity
rules then establish equivalence of the corresponding
atomic membership or identity formulas.

Thus, realization invariance holds for atomic formulas.

The logical connectives preserve the equivalence
by ordinary propositional reasoning.

At quantifiers, the bound variable is renamed
alphabetically whenever necessary to prevent capture.
The induction hypothesis then yields equivalence
within the quantifier's scope, and Quantifier
Congruence establishes equivalence of the quantified
formulas.

Defined notation is handled through its prescribed
expansion. Consequently, the equivalence holds for
all admitted extended formulas.

We have therefore established
\[
\{
\Real_T(\eta)
\}
\vdash
\varphi[T/\xi]
\leftrightarrow
\varphi[\eta/\xi].
\]

To obtain the second assertion, recall that
\[
\vdash
\exists\eta\,\Real_T(\eta),
\]
since \(\eta\) was chosen fresh for \(T\).

Suppose first that \(\varphi[T/\xi]\) holds.
Existential elimination supplies a realization
\(\eta\) of \(T\). By the equivalence just proved,
\(\varphi[\eta/\xi]\) follows. Hence
\[
\exists\eta
\bigl(
\Real_T(\eta)
\land
\varphi[\eta/\xi]
\bigr).
\]

Conversely, suppose that this existential formula
holds. Existential elimination supplies a variable
\(\eta\) satisfying both
\[
\Real_T(\eta)
\qquad\text{and}\qquad
\varphi[\eta/\xi].
\]
Realization invariance gives
\(\varphi[T/\xi]\).

The eigenvariable conditions are satisfied because
\(\eta\) was chosen fresh for \(\varphi\) and \(T\).

The two implications establish the required
equivalence.
\end{proof}

Realization invariance is the central result needed
for the ordinary term-like use of class-builder
notation. It ensures that substitution of a class
term can always be replaced by reasoning with a
fresh variable realizing that term.

In particular, any derivation using such a substitution
can be reduced to a derivation using only variables,
realizations, and the original inference rules.


\nbgpara{5}{Admissibility of Extended-Term Rules}

We can now establish that the familiar inference rules
which involve terms remain admissible when their
variable arguments are replaced by arbitrary extended
terms.

\begin{metatheorem}[Admissibility of Extended-Term Rules]
\label{meta:extended-term-rules}
Let \(T\) be an admitted extended term, let \(\xi\)
be a variable, and let \(\varphi\) be an extended
formula.

Whenever the indicated substitution is legitimate,
the following rules are admissible:
\[
\frac{\forall\xi\,\varphi}
{\varphi[T/\xi]}
\quad
\forall\text{-E},
\qquad
\frac{\varphi[T/\xi]}
{\exists\xi\,\varphi}
\quad
\exists\text{-I}.
\]

Identity introduction is admissible for arbitrary
extended terms:
\[
\frac{}{T=T}
\quad
=\text{-I}.
\]

Moreover, identity elimination extends to arbitrary
extended terms. In particular, from
\[
S=T
\qquad\text{and}\qquad
\varphi[S/\xi]
\]
one may infer
\[
\varphi[T/\xi],
\]
provided the substitutions are legitimate.

The previously established derived rules involving
sethood, subclasshood, and membership likewise
admit extended-term instances wherever their
formation conditions are satisfied.
\end{metatheorem}

\begin{proof}
For variable terms, the assertions are instances
of the original rules or their previously established
derived forms.

Suppose therefore that \(T\) is a class term.

For universal elimination, choose a fresh variable
\(\eta\). By the Class-Term Extension Metatheorem,
\[
\vdash
\exists\eta\,\Real_T(\eta).
\]

Under the assumption \(\Real_T(\eta)\), the primitive
universal-elimination rule permits the inference
\[
\forall\xi\,\varphi
\quad\Longrightarrow\quad
\varphi[\eta/\xi].
\]

By Realization Invariance,
\[
\varphi[\eta/\xi]
\quad\Longrightarrow\quad
\varphi[T/\xi]
\]
under the same assumption.

Existential elimination then removes the realization
assumption and yields the required extended-term
instance of universal elimination.

For existential introduction, suppose that
\[
\varphi[T/\xi]
\]
has been derived.

Using the existence of a realization and Realization
Invariance, we obtain
\[
\exists\eta
\bigl(
\Real_T(\eta)
\land
\varphi[\eta/\xi]
\bigr).
\]

Within the corresponding existential-elimination
subderivation, primitive existential introduction
gives
\[
\exists\xi\,\varphi.
\]
The realization variable can then be discharged,
yielding the required rule.

For identity introduction, choose a fresh realization
variable \(\eta\). From \(\Real_T(\eta)\), identity
reflexivity gives \(\eta=\eta\). The representation
equivalence for atomic identity formulas then yields
\(T=T\). Existence of a realization permits the
assumption to be discharged.

Finally, consider identity elimination.

Choose fresh representatives \(\alpha\) and \(\beta\)
of \(S\) and \(T\). Under their representation
conditions, the premise
\[
S=T
\]
is equivalent to
\[
\alpha=\beta.
\]

By Realization Invariance, the other premise
\[
\varphi[S/\xi]
\]
is equivalent, under the corresponding representation
condition, to
\[
\varphi[\alpha/\xi].
\]

The primitive identity-elimination rule, applied
through primitive expansions, therefore yields
\[
\varphi[\beta/\xi].
\]

A further application of Realization Invariance
gives
\[
\varphi[T/\xi].
\]

The representative variables are then eliminated
using the existence and uniqueness of representatives.

The remaining derived rules are treated in the
same manner. Each extended-term application is
replaced by reasoning with representatives, followed
by the corresponding variable-based rule.

Thus, no new primitive inference rules are required.
\end{proof}

In particular, the derived Set-I rule now permits
the inference
\[
\frac{S\in T}{\Set(S)}
\quad
\mathrm{Set}\text{-I}
\]
for arbitrary admitted extended terms \(S\) and \(T\).

Likewise, the previously established rules for
subclasshood may be applied to extended terms.

The fundamental Membership Law also admits
extended-term arguments.

\begin{metatheorem}[Generalized Membership and Identity Laws]
\label{meta:generalized-class-term-laws}
Let
\[
T=\{\xi\mid\varphi\}
\]
be an admitted general class term, and let \(S\)
be an admitted extended term.

Whenever the indicated substitution is legitimate,
\[
\boxed{
\vdash
S\in T
\leftrightarrow
\bigl(
\Set(S)\land\varphi[S/\xi]
\bigr).
}
\]

Moreover, if \(\alpha\) is chosen fresh for \(S\)
and \(T\), then
\[
\boxed{
\vdash
S=T
\leftrightarrow
\forall\alpha
\bigl(
\alpha\in S
\leftrightarrow
\operatorname{Cond}_T(\alpha)
\bigr).
}
\]
\end{metatheorem}

\begin{proof}
For membership, choose a variable \(\zeta\) fresh
for \(S\) and \(T\).

The variable-argument Membership Law gives
\[
\vdash
\zeta\in T
\leftrightarrow
\bigl(
\Set(\zeta)
\land
\varphi[\zeta/\xi]
\bigr).
\]

Since \(\zeta\) is fresh for the parameters of this
formula, universal introduction and the admissible
extended-term universal-elimination rule permit
its instantiation with \(S\).

Capture-avoiding substitution and the previously
established compatibility properties then give
\[
\vdash
S\in T
\leftrightarrow
\bigl(
\Set(S)\land\varphi[S/\xi]
\bigr).
\]

For identity, the atomic equality reduction gives
\[
\vdash
S=T
\leftrightarrow
\forall_S\alpha
\bigl(
M_S(\alpha)\leftrightarrow
M_T(\alpha)
\bigr).
\]

By the variable-argument Membership Law,
\(M_S(\alpha)\) is equivalent to \(\alpha\in S\),
and \(M_T(\alpha)\) is equivalent to
\(\operatorname{Cond}_T(\alpha)\).

Both membership conditions imply sethood.
Consequently, the restriction of the universal
quantifier may be removed without changing
the equivalence.

This gives
\[
\vdash
S=T
\leftrightarrow
\forall\alpha
\bigl(
\alpha\in S
\leftrightarrow
\operatorname{Cond}_T(\alpha)
\bigr),
\]
as required.
\end{proof}

Notice that the generalized Identity Law does not
require us to define \(\Real_T(S)\) for arbitrary
extended terms. Identity with a class term is
expressed directly through the corresponding
extensional membership condition.

The four class-builder rules established previously
may now be used with extended-term arguments through
these generalized laws.


\nbgpara{6}{Generalized Bounded Class Terms}

We can finally extend bounded class-builder notation
from variable bounds to arbitrary extended-term bounds.

\begin{definition}[Generalized Bounded Class Terms]
\label{def:generalized-bounded-class-terms}
Let \(S\) be an admitted extended term, let \(\xi\)
be a variable not occurring free in \(S\), and let
\(\varphi\) be a predicative extended formula.

Then
\[
\{\xi\in S\mid\varphi\}
\]
is an admitted bounded class term, provided that
its defining condition is formed from previously
admitted expressions.

It is defined by the expansion
\[
\{\xi\in S\mid\varphi\}
\defequiv
\{\xi\mid(\xi\in S\land\varphi)\}.
\]
\end{definition}

The condition that \(\xi\) not occur free in \(S\)
ensures that the builder does not inadvertently bind
free variables belonging to its bounding term.

The builder rank is inherited from the general
class term to which the bounded term is expanded:
\[
\operatorname{br}
\bigl(
\{\xi\in S\mid\varphi\}
\bigr)
=
1+
\max
\bigl(
\operatorname{br}(S),
\operatorname{br}(\varphi)
\bigr).
\]

Since the atomic formula
\[
\xi\in S
\]
has a justified predicative reduction, the condition
\[
\xi\in S\land\varphi
\]
is predicative whenever \(\varphi\) is predicative.

Thus, generalized bounded class terms are admitted
by the existing class-term construction. They require
no new class-forming axiom and no independent
existence theorem.

Their fundamental laws follow immediately from
the generalized laws for general class terms.

\begin{metatheorem}[Generalized Bounded Class-Builder Laws]
\label{meta:generalized-bounded-builder-laws}
Let
\[
T=\{\xi\in S\mid\varphi\}
\]
be an admitted generalized bounded class term,
and let \(U\) be an admitted extended term.

Whenever the indicated substitution is legitimate,
\[
\vdash
U\in T
\leftrightarrow
\bigl(
U\in S
\land
\varphi[U/\xi]
\bigr).
\]

Moreover, for a variable \(\alpha\) chosen fresh
for the displayed terms and formulas,
\[
\vdash
U=T
\leftrightarrow
\forall\alpha
\left(
\alpha\in U
\leftrightarrow
\bigl(
\alpha\in S
\land
\varphi[\alpha/\xi]
\bigr)
\right).
\]
\end{metatheorem}

\begin{proof}
Expand the bounded term according to its definition:
\[
T
\defequiv
\{\xi\mid(\xi\in S\land\varphi)\}.
\]

The generalized Membership Law gives
\[
\vdash
U\in T
\leftrightarrow
\left(
\Set(U)\land
\bigl(
U\in S
\land
\varphi[U/\xi]
\bigr)
\right).
\]

Since the extended-term Set-I rule gives
\[
\vdash
U\in S\to\Set(U),
\]
the sethood conjunct is redundant. This proves
the first equivalence.

The second follows from the generalized Identity
Law by the same simplification, applied within
the unrestricted membership specification.
\end{proof}

Consequently, expressions such as
\[
\{\xi\in V\mid\varphi\}
\]
are legitimate whenever \(V\) is an admitted
extended term and the stated predicativity and
variable-binding conditions are satisfied.

This completes the justification of bounded
class-builder notation in its general form.


\nbgpara{7}{Eliminability and Conservativity}

We conclude by making explicit the relation between
the extended syntax and the original deductive system.

\begin{metatheorem}[Conservativity of Class-Term Notation]
\label{meta:conservativity-of-class-term-notation}
Every extended formula containing admitted class terms
has a primitive expansion, determined up to alphabetic
variation.

Every derivation using class terms, their defined
relations, and the derived extended-term rules
abbreviates a derivation using formulas of the
primitive language and the original inference rules
and axioms of NBG.

Consequently, the extension by class-builder notation
is conservative: it proves no new formula of the
primitive language that was not already derivable
from the original axioms and inference rules.
\end{metatheorem}

\begin{proof}
The existence and coherence of primitive expansions
were established by the Class-Term Extension Metatheorem.

The fundamental Membership and Identity Laws were
derived using Class Formation, Extensionality, and
the original natural deduction rules.

The predicative reductions assigned to atomic
formulas containing class terms were likewise
proved equivalent to those formulas. Their use
inside class-forming conditions therefore supplies
only derived instances of the original Class
Formation Schema.

Realization Invariance establishes that substitution
of an extended term may be replaced by reasoning with
a fresh variable representing that term. The
extended-term inference rules were obtained by
applying the primitive rules to such representatives
and subsequently eliminating the realization
variables.

Finally, the generalized bounded class terms and
the extended forms of previously defined relations
are eliminable through their prescribed expansions.

It follows by induction on the construction of an
extended derivation that every occurrence of an
extended formula may be replaced by a primitive
expansion and every extended inference step by a
finite derivation using the original rules.

Where different legitimate choices of fresh bound
variables yield alphabetically variant expansions,
the corresponding replacement is licensed by the
previously established metatheorems concerning
alphabetic variants.

Thus, any derivation of a primitive formula using
the extended syntax can be transformed into a
derivation of that same formula in the original
deductive system.

The extension therefore introduces no new
deductive power.
\end{proof}

The significance of this result is that class-builder
notation may henceforth be used with the ordinary
freedom of mathematical term notation, subject to
the stated formation and predicativity conditions.

In particular, expressions such as
\[
\Set(R),\qquad
\PC(R),\qquad
A\subseteq V,\qquad
R\in R,
\]
and
\[
\{x\in V\mid\varphi\}
\]
are formally legitimate whenever their constituent
terms and defining conditions are admitted.

Likewise, the ordinary quantifier and identity rules
may be applied to extended terms through the derived
rules established above.

None of these conveniences changes the primitive
language of NBG. Variables remain its only terms,
and \(\in\) remains its only non-logical primitive
relation symbol.

Class-builder expressions and their accompanying
rules are therefore entirely eliminable. They
provide a more convenient notation for reasoning
within NBG, rather than an enlargement of the
underlying theory.